% =====================================================================
%  MIC Summer-of-Code  --  Weeks 6.5 to 8
%  Paper Reading + Guided Implementation  ->  Final Project
%
%  PLACEHOLDER NOTICE.  Compile with pdflatex (once is enough) or upload
%  to Overleaf.  The mentor replaces / supplements this with the actual
%  paper to implement.
%      pdflatex Paper_To_Implement_PLACEHOLDER.tex
% =====================================================================
\documentclass[11pt]{article}

\usepackage[a4paper,margin=1in]{geometry}
\usepackage{enumitem}
\usepackage{xcolor}
\usepackage{microtype}
\usepackage[most]{tcolorbox}
\usepackage{titlesec}
\usepackage[colorlinks=true,linkcolor=blue!45!black,urlcolor=blue!55!black]{hyperref}

\titleformat{\section}{\large\bfseries\sffamily\color{blue!40!black}}{\thesection}{0.5em}{}

\newtcolorbox{keybox}[1][]{colback=blue!4!white,colframe=blue!40!black,
  fonttitle=\bfseries\sffamily,title=#1,boxrule=0.7pt,arc=2pt,left=5pt,right=5pt,top=4pt,bottom=4pt}
\newtcolorbox{waitbox}[1][]{colback=orange!8!white,colframe=orange!60!black,
  fonttitle=\bfseries\sffamily,title=#1,boxrule=0.9pt,arc=2pt,left=5pt,right=5pt,top=4pt,bottom=4pt}

\title{\textbf{Paper to Implement}\\[2pt]
\large Weeks 6.5--8: Paper Reading, Guided Implementation \& Final Project\\[6pt]
\normalsize Medical Image Computing -- Summer of Code}
\author{}
\date{}

\begin{document}
\maketitle
\vspace{-2.2em}

\begin{waitbox}[$\boldsymbol{\Large\,\text{To be updated --- please wait}}$]
\large
The specific \textbf{paper you will implement} for this block (together with any
data and starter code) will be \textbf{uploaded here by your mentor}.

\smallskip
Please \textbf{wait for that upload} before starting the implementation. This page
is a placeholder; check back, or ask your mentor when the paper will be posted.
\end{waitbox}

\bigskip

% =====================================================================
\section{What to do while you wait}
% =====================================================================

You do not need to sit idle. Get everything ready so that, the moment the paper
lands, you can start implementing:

\begin{itemize}[leftmargin=*,nosep]
  \item \textbf{Set up the environment.} Make sure the \texttt{mic\_soc}
        environment runs (NumPy / SciPy / Matplotlib, and PyTorch + MONAI for the
        deep-learning weeks). Confirm you can open a notebook and load a dataset.
  \item \textbf{Refresh the paper-reading method.} Re-read the \emph{three-pass}
        approach in \texttt{week\_7/Readme.md} (\S2) --- you will use it on the
        paper your mentor provides.
  \item \textbf{Skim the starter pool.} Browse the example papers and datasets in
        \texttt{week\_7/Readme.md} (\S4--5) and the project rubric in
        \texttt{week\_8/Readme.md}, so you know what ``good'' looks like.
  \item \textbf{Be able to fetch data.} Check you can download a small dataset
        (e.g.\ MedMNIST) and run one training step end-to-end on five examples.
\end{itemize}

% =====================================================================
\section{What you will be asked to do (once the paper is posted)}
% =====================================================================

\begin{enumerate}[leftmargin=*,nosep]
  \item \textbf{Read the paper} deeply (three passes) and write a one-page summary
        (use the template in \texttt{week\_7/Readme.md}).
  \item \textbf{Reproduce one experiment} --- one figure or one number --- from the
        paper, on a manageable dataset.
  \item \textbf{Run a baseline} (a classical method or a simple network) on the
        same data, to give the result context.
  \item \textbf{Write it up}: what the paper claimed, what you reproduced, the gap,
        and what it taught you --- then fold it into the \textbf{final project}
        (\texttt{week\_8/Readme.md}: report + slides + reproducible code).
\end{enumerate}

\begin{keybox}[Where things live]
Full week plans, the paper-reading template, the starter paper pool, datasets,
and the final-project rubric are in \texttt{week\_7/Readme.md} and
\texttt{week\_8/Readme.md}. This document only marks the \emph{paper to implement}
as pending.
\end{keybox}

\vfill
\begin{center}\small\itshape
MIC Summer-of-Code. Placeholder notice --- to be replaced/supplemented by the
mentor with the actual paper to implement.
\end{center}

\end{document}
